% ============================================================
%  ch3_enhanced  三、幂零锥尖端处的增强镜像对称（原书 pp. 2235--2246）
% ============================================================

\section{幂零锥尖端处的增强镜像对称}\label{enhanced}

\origpage{2235}下面这些考察的原初动机，是要找到一种办法来检验 \cite{baraglia-schaposnik,garcia-prada-etal,hitchin-characteristic} 所提出的 $(B,A,A)$ 型与 $(B,B,B)$ 型膜的猜想镜像对。迄今为止唯一的检验——常常就在这些文献中完成——是检查所提议的镜像对相对于 Hitchin 映射是否确实互为 Fourier--Mukai 对偶。而这种检验只能在 Hitchin 底上的\emph{一般}点处进行。我们希望看到更多整体性的检验，特别是那些能在 Hitchin 映射的 $0$-纤维，即整体幂零锥上验证镜像对称提议的检验。

我们在 \cite{hausel-mellit-pei,hausel-hitchin} 中引入了一种技术：考察镜像对称对相应范畴中态射的作用，从而在幂零锥上验证镜像对称的提议。在我们的范畴中考察态射的困难，来自模空间的非紧性。例如导出范畴 $D(\M_\Dol)$ 中态射构成的向量空间通常是无穷维的。为了度量它们的大小，我们要看其上的 $\T$-等变结构。回忆一下：复数乘法群 $\T:=\C^\times$ 通过 $\lambda:(E,\Phi)\to (E,\lambda\Phi)$（即把 Higgs 场作数乘）作用在 $\M$ 上。这里我们感兴趣的是几何 Langlands 对应的经典极限 \eqref{classlim} 的一个 $\T$-等变推广，作为一阶近似，它应当是 $\M_\Dol$ 与 $\M^\vee_{\Dol}$ 的 $\T$-等变导出范畴之间的一个等价
\[
\mathscr{S}:D_\T(\M_\Dol) \sim D_\T(\M^\vee_{\Dol}).
\]
$D_\T(\M_\Dol)$ 中两个对象 $\calF_1$、$\calF_2$ 之间的态射，可以等同于分次向量空间 $\Ext^*(X;\calF_1,\calF_2)$。为了度量这个分次向量空间，我们注意到 $\T$ 作用在它上面；再假定各权空间有限维、且当权足够大时消失（对半射影簇我们正是这样预期的），就可以定义\textbf{等变 Euler 型}为
\begin{multline*}
\chi_\T(X;\calF_1,\calF_2)=\sum_{k,l} \dim\bigl(H^k({R}{\mathcal Hom}(\calF_1,\calF_2))^l\bigr) (-1)^k t^{-l}=\\
=\sum_{k,l} \dim\Bigl(\Hom_{{\mathbf D}_{coh}(X)}(\calF_1,\calF_2[k])^l\Bigr) (-1)^k t^{-l}=\\
=\sum_{k,l} \dim\bigl(\Ext^k(X;\calF_1,\calF_2)^l\bigr) (-1)^k t^{-l} \in \C((t)).
\end{multline*}
有了它，我们期望 $\mathscr{S}$ 是一个等距：
\[
\chi_\T(\mathscr{S}(\calF),\mathscr{S}(\calG))=\chi_\T(\calF,\calG).
\]

在 \cite{hausel-mellit-pei} 中，我们对 \cite{hitchin-characteristic} 中若干对猜想镜像膜验证了这一等距；而在 \cite{hausel-hitchin} 中，我们验证了与本文有关的那些猜想镜像对之间的这一等距，详见下文。在本文第二部分，我们将回顾 \cite{hausel-hitchin} 中关于非常稳定上行流的镜像的一些结果，接着依照 \cite{hausel-hitchin1} 引入并研究拉格朗日上行流的重数代数，最后考察重数代数在镜像上应当对应什么。

\subsection{非常稳定 Higgs 丛与镜像对称}

出发点是新近的论文 \cite{hausel-hitchin}。先回顾它的框架。

\subsubsection{半射影簇的 Białynicki--Birula 分解}

\origpage{2236}设复数乘法群 $\T:=\C^\times$ 作用在一个簇（即 $\C$ 上有限型的既约分离概形，可以可约）上。称此作用是\textbf{半射影的} \cite[\S 1]{hausel-large}，如果下面三条成立：

\begin{enumerate}
	\item 作用是\emph{线性的}，即存在 $X$ 到 $\P^N$ 的局部闭 $\T$-等变嵌入，且 $\T$ 在 $\P^N$ 上的作用线性（例如 $X$ 正规且拟射影时）；
	\item 不动点子簇 $X^\T$ 是 proper 的，从而是射影的；
	\item 对一切 $x\in X$，$\lim_{\lambda\to 0} \lambda\cdot x$ 存在。
\end{enumerate}
对 $\alpha\in X^\T$，定义 $W^+_\alpha:=\{x\in X\mid \lim_{\lambda\to 0} \lambda\cdot x = \alpha\}$，称为自 $\alpha$ 出发的\textbf{上行流}；并定义 $W^-_\alpha:=\{x\in X\mid \lim_{\lambda\to \infty} \lambda\cdot x = \alpha\}$，称为自 $\alpha$ 出发的\textbf{下行流}。于是有 $X$ 的 \textbf{Białynicki--Birula 分解}
\[
X=\coprod_{\alpha\in X^\T} W_\alpha^+,
\]
并定义射影簇
\[
\calC:=\coprod_{\alpha\in X^\T} W_\alpha^-
\]
为 $X$ 的\textbf{芯}。由此得到

\begin{theorem}[{\cite{bialynicki-birula},\cite[Proposition 2.1, Proposition 2.10]{hausel-hitchin}}]
当 $\alpha\in (X^s)^\T$ 是光滑轨迹上的 $\T$-不动点时，$W^{\pm}_\alpha\subset X$ 是局部闭子簇，且作为 $\T$-簇有 $W^{\pm}_\alpha\cong T^{\pm}_\alpha X$。进而，当 $\omega\in \Omega^2(X^s)$ 是齐次度为 $1$ 的辛形式时，$W^+_\alpha\subset X$ 与 $\calC\subset X$ 都是拉格朗日子簇。
\end{theorem}

当 $W^+_\alpha\cap \calC =\{\alpha\}$ 时，我们称 $\alpha\in (X^s)^\T$ 与 $W^+_\alpha$ 是\textbf{非常稳定的}。等价地说，$\alpha\in (X^s)^\T$ 非常稳定，当且仅当 $W^+_\alpha\cap W_\beta^-\neq \emptyset$ 蕴含 $\beta=\alpha$。更一般地可以证明：当 $W^+_\alpha\cap W_\beta^-\neq \emptyset$ 时规定 $\alpha\leq \beta$，由此得到的 $\leq$ 是一个偏序；而 $\alpha\in (X^s)^\T$ 非常稳定，当且仅当它关于这个序是极大的。于是有

\begin{theorem}[{\cite[Proposition 2.14]{hausel-hitchin}}]\label{closed}
$\alpha\in (X^s)^\T$ 非常稳定，当且仅当 $W_\alpha^+\subset X$ 是闭的。
\end{theorem}

\subsubsection{Higgs 丛的 Białynicki--Birula 分解}

以下取 $\G=\GL_n$。用 $\M$ 记亏格为 $g$ 的光滑射影曲线 $C$ 上秩 $n$、次数 $d$ 的半稳定 Higgs 丛 $(E,\Phi)$ 的模空间。这里 $E$ 是 $C$ 上秩 $n$、次数 $d$ 的向量丛，Higgs 场 $\Phi\in H^0(C;\End(E)\otimes K)$。

由 Higgs 场的特征多项式给出 \textbf{Hitchin 映射} \cite{hitchin-stable}：
\[
\begin{array}{cccc}
h: &\M&\to& \calA:=\times_{i=1}^n H^0(C;K^i)\\
 &(E,\Phi)&\mapsto& \det(x-\Phi).
\end{array}
\]
于是 $h$ 是 proper 映射 \cite{hitchin,nitsure,simpson}，并关于 $\M$ 上一个自然的全纯辛形式构成完全可积的 Hamilton 系统 \cite{hitchin-stable,faltings}。特别地，一般纤维是拉格朗日的阿贝尔簇，即某些谱曲线的 Jacobi 簇。

$\T$ 在 $\M$ 上的作用由 $(E,\Phi)\mapsto (E,\lambda\Phi)$ 给出。若让 $\T$ 以权 $i$ 作用在 $H^0(C;K^i)$ 上，则 Hitchin 映射是 $\T$-等变的。由于这些权在 $\calA\cong \times_{i=1}^n H^0(C;K^i)$ 上全为正，该作用（在 $\calA$ 上）是半射影的；又因 Hitchin 映射 proper 且 $\T$-等变，可知 $\T$ 在 $\M$ 上的作用也是半射影的。此外还可得到：$\calA$ 的芯，即原点 $0\in \calA$，其原像正是 $\M$ 的芯。对 $\M$ 而言，芯与幂零锥 $h^{-1}(0)_{red}=\calC$ 一致。

把 Drinfeld 与 Laumon \cite{drinfeld2,laumon1} 的非常稳定丛的概念推广，我们于是可以定义\textbf{非常稳定 Higgs 丛}：它是稳定的 $\T$-不动 Higgs 丛 $\calE\in {\M^s}^\T$，并且满足 $W^+_\calE\cap h^{-1}(0)=\{\calE\}$，即它的上行流中唯一的幂零 Higgs 丛就是它自己。于是由定理 \ref{closed}，$\calE\in {\M^s}^\T$ 非常稳定，恰好就是它的上行流为闭集。为了用 Hitchin 映射重新表述，我们有定理 \ref{closed} 的另一个版本。

\begin{theorem}[\cite{hausel-hitchin1}，参看 \cite{zelaci}]\label{proper}
$\calE \in {\M^s}^\T$ 非常稳定，当且仅当 $h^{-1}({0})\cap W^+_\calE=\{\calE\}$；当且仅当 $h_\calE:=W_\calE^+\to \calA$ 是 proper 的；当且仅当它是有限的。
\end{theorem}

\begin{definition}\label{equmult}
对 $\calE\in \M^{s\T}$，定义有理函数
\[
m_{\calE}(t):=\frac{\chi_\T(\Sym(T^{+*}_\calE))}{\chi_\T(\Sym(\calA^*))}\in \Z(t).
\]
我们称之为 $W^+_\calE$ 的\textbf{等变重数}。
\end{definition}

于是有

\begin{theorem}[{\cite[Corollary 5.4]{hausel-hitchin}}]\label{eqmul}
当 $\calE\in \M^{s\T}$ 非常稳定时，$m_{\calE}(t)$ 是一个多项式：\begin{enumerate}\item 系数非负；\item 是回文的（palindromic）；\item 是首一的；\item 满足 $m_{\calE}(1)=m_{F_\calE}$，即幂零锥中分量 $N_{F_\calE}$ 的重数。\end{enumerate}
\end{theorem}

设 $\ell\in \Z$、$m_i\in \N$，并设
\[
\delta:=(\delta_0,\delta_1,\dots,\delta_{n-1})\in Jac_\ell(C)\times C^{[m_1]}\times \dots \times C^{[m_{n-1}]}
\]
是 $C$ 上一列代表除子。由此可以构造一个 $(1,\dots,1)$ 型的 $\T$-不动 Higgs 丛 $\calE_\delta=(E_\delta,\Phi_\delta)$：其中
\[
\calE_\delta=M_0\oplus\dots \oplus M_{n-1}
\]
是秩 $n$ 的向量丛，
\[
M_i:=\calO(\delta_0+\dots+\delta_i)K^{-i},
\]
而
\[
\Phi_\delta|_{M_i}:M_i\to M_{i+1}K\subset E_\delta K
\]
由
\[
H^0(C;M_i^{-1}M_{i+1}K)\cong H^0(C;\calO(\delta_i))
\]
的定义截面给出。

下面这个定理分类了全部非常稳定的 $\T$-不动 $(1,\dots,1)$ 型 Higgs 丛，并给出它们的等变重数。

\begin{theorem}[{\cite[Theorem 4.16, (5.18)]{hausel-hitchin}}]\label{mult111}
设 $\delta$ 如上，并设 $\calE_\delta$ 是稳定 Higgs 丛。则 $\calE_\delta$ 非常稳定，当且仅当有效除子 $\delta_1+\dots+\delta_{n-1}$ 是既约的。其等变重数为
\[
m_{\calE_\delta}(t)=\prod_{i=1}^{n-1}\left[\begin{array}{c} n \\ i \end{array}\right]^{m_i}_t,
\]
即 $t$-二项式系数的乘积。
\end{theorem}

\subsection{重数代数与拉格朗日子流形上的显式 Hitchin 系统}

\origpage{2238}主要想法是：对 $\calE\in \M^{s\T}$，在 Arnold 学派 \cite{arnold-etal} 的框架中研究受限 Hitchin 映射

\[
h_\calE:=h|_{W^+_\calE}: W^+_\calE\to \calA.
\]
这是同维的半射影向量空间（即只有正的 $\T$-权）之间的一个 $\T$-等变\textbf{拉格朗日映射}。\cite[\S 12.3]{arnold-etal} 称这类映射为\textbf{拟齐次的}。

由定理 \ref{proper} 回忆：$\calE$ 非常稳定，当且仅当 $h_\calE^{-1}(0)=\{\calE\}$，当且仅当 $h_\calE$ 是 proper 的。对这类映射——在\emph{同上}文献中称为\textbf{非退化的}——\cite[\S 4,5,12]{arnold-etal} 引入并研究了它的局部重数代数。

\begin{definition}
当 $\calE\in \M^{s\T}$ 非常稳定时，定义
\[
Q_\calE:=Q_{h_\calE}:= \C[W^+_\calE]/(h_\calE^{-1}({\mathfrak m}_0))=\C[W^+_\calE]/(h_1,\dots,h_N),
\]
称之为 $h_\calE$ 在 $\calE$ 处的\textbf{局部重数代数}。这里 ${\mathfrak m}_0\subset \C[\calA]$ 是 $0\in \calA$ 处的极大理想，而在某组齐次坐标下
\[
h_\calE=(h_1,\dots,h_N):\C^N\cong W^+_\calE\to \C^N\cong \calA.
\]
\end{definition}

从概形的角度看，$Q_\calE$ 就是 $h_\calE$ 在 $0$ 上的概形论纤维的坐标环，也就是上行流与幂零锥的概形论交 $W^+_\calE\cap h^{-1}(0)$ 的坐标环。因为 $h_\calE$ 是 $\T$-等变的，所以 $Q_\calE$ 上带有 $\T$-作用，从而带有分次结构。由于这一点，有时我们也把 $Q_\calE$ 称为 $h_\calE$ 在 ${\calE}$ 处的\textbf{等变重数代数}。

\begin{remark}
注意：用 $N$ 个生成元与 $N$ 个关系把代数 $Q_\calE$ 显式地确定出来，就等于在 $W^+_\calE$ 上给出了一组坐标，使得 Hitchin 映射\emph{显式地}由这些关系给出。
\end{remark}

利用 \cite[\S 4,5,12]{arnold-etal} 的结果，我们有

\begin{theorem}[\cite{hausel-hitchin1}]
设 $\calE\in \M^s$ 非常稳定。则它的局部重数代数具有下述性质：\begin{enumerate}\item 有限维；\item 分次，$Q_\calE:=\bigoplus_{k=0}^m Q_\calE^k$，且 $Q_\calE^0\cong \C$；\item 是 Gorenstein 的，其基座（socle）$Q^m_\calE\cong \C J_{h_\calE}$ 是一维的，由 $h_\calE=(h_1,\dots,h_N)$ 的 Jacobian $J_{h_\calE}$ 张成；\item 是 Poincaré 对偶环：存在自然的双线性配对 $(,):Q_\calE\times Q_\calE \to \C$，对一切 $k$ 都诱导出完美配对 $Q_\calE^k\times Q_\calE^{m-k}\to \C$；\item 最后，它的 Poincaré 多项式 $\sum_{k=0}^m \dim(Q^k_\calE) t^k=m_\calE(t)$ 与等变重数一致。\end{enumerate}
\end{theorem}

这个结果对我们在定理 \ref{eqmul} 中观察到的等变重数多项式的所有性质给出了圆满的解释。下面这个结果则解释了量子二项式系数——众所周知它们是 Grassmann 簇的 Poincaré 多项式——为何会出现在等变重数之中。

\begin{theorem}[\cite{hausel-hitchin1}]\label{mainema}
设 $\delta=(\delta_0,\delta_1,\dots,\delta_{n-1})\in J_\ell(C)\times C^{[m_1]}\times \dots \times C^{[m_{n-1}]}$ 是一列除子。假定 $\delta_1+\dots+\delta_{n-1}$ 既约，则由定理 \ref{mult111} 知 $\calE_\delta$ 非常稳定，并且它的等变重数代数为
\[
Q_{\delta}\cong \prod_{i=1}^{n-1} H^*(\Gr(i,n);\C)^{m_i},
\]
即 $\C^n$ 中 $i$ 维子空间的 Grassmann 簇 $\Gr(i,n)$ 的复上同调环之积。
\end{theorem}

我们对定理 \ref{mainema} 的攻破，始于用 Hecke 变换来理解映射 $h_{\delta}:W^+_{\delta}\to \calA$。这里先解释第一个非平凡的情形：
\[
\delta^k_c:=(\delta_0,\dots,\delta_{n-1}),
\]
其中除 $i=k$ 之外 $\delta_i=0$，而除子 $\delta_k=c$ 只含一个点。为此我们考虑 Hitchin 截面 $W^+_0:=W^+_{\calE_0}$ 上的第 $k$ 个基本 Hecke 对应；这里 $W^+_0$ 与\textbf{标准均匀化 Higgs 丛} $\calE_0=(E_0,\Phi_0)$ 的上行流一致。回忆 \cite[(3.16)]{hausel-hitchin}，用链式记号它写成
\[
(E_0,\Phi_0)=(\calO\oplus K^{-1}\oplus \dots\oplus K^{1-n},\Phi_0)= \calO\stackrel{1}{\rightarrow} K^{-1} \stackrel{1}{\rightarrow} \dots \stackrel{1}{\rightarrow}  K^{1-n},
\]
其中箭头上的态射都扭上 $K$。标准均匀化 Higgs 丛位于幂零锥的顶端。我们将用 Hecke 变换，从 Hitchin 截面 $W^+_0$ 生成其它一切 $(1,\dots,1)$ 型上行流。

令
\begin{equation}\label{fundamental}
\calH_k:=\left\{ (\calE_a,V)\in W^+_0\times \Gr(k, E_0|_{c})\;\middle|\; \Phi_a|_{c}(V)\subset V\right\}\subset W^+_0\times \Gr(k, E_0|_{c}).
\end{equation}
它可以在 $\calA$ 上 étale 局部地构造出来，作为 $\Gr(k,E_0|_{c})$ 的一个自映射的不动点概形；该自映射由可逆的 $\Phi_a|_{c}-\lambda I :E_0\to E_0$ 诱导，其中 $\lambda\in \C$ 不是 $\Phi_a$ 的特征值。对 $(\calE_a,V)\in \calH_k$，在 $V$ 处对 $\calE_a$ 作 Hecke 变换，就得到 $\calH_k\cong W^+_{\delta^k_c}$；并且有交换图
\[
\begin{array}{ccc}
\calH_k&\stackrel{\cong}{\longrightarrow}& W^+_{\delta^k_c}\\
\pi_k{\Big\downarrow} && {\Big\downarrow} h_{{\delta^k_c}} \\
W^+_0 &\stackrel{h_0}{\longrightarrow}& \calA
\end{array}
\]
因此 $Q_{\delta^k_c}\cong Q_{\pi_k}$。而 $Q_{\pi_k}$ 的计算可以在 Grassmann 簇 $\Gr(k, E_0|_{c})$ 中完成，并用两种不同的方式给出下面两个同构：
\begin{multline}\label{cohgras}
H^*(\Gr(k, E_0|_{c});\C)\cong Q_{\pi_k}\cong  \\
\cong \frac{\C[p_1,\dots,p_k,q_1,\dots,q_{n-k}]}{\left((p_k+\dots+p_1x^{k-1}+x^k)(q_{n-k}+\dots+q_1x^{n-k-1}+x^{n-k})=x^n\right)}.
\end{multline}
当然，第一个环与第三个环同构，给出的正是 Grassmann 簇上同调环的著名表述。

事实上，我们可以把 \eqref{cohgras} 中重数代数的确定，看成是在上行流 $W^+_{\delta^k_c}$ 上给出坐标，使得 Hitchin 映射 $h_{{\delta^k_c}}$ 变得显式——基本上就由 Grassmann 簇上同调环中的那些关系给出。

\subsection{wobbly 拉格朗日子流形上的显式 Hitchin 系统}

\origpage{2241}这里我们说明，如何用仿射 Grassmannian \cite{gortz,zhu} 把上述技术推广，以理解 wobbly 的——即不非常稳定的——上行流的重数代数。\footnote{译注：“wobbly”是原文自造的说法，指不非常稳定的上行流；原文在此没有给出更细的定义，其含义只能由定理 \ref{mainema} 与定理 \ref{conjecturewobbly} 的对照看出：非常稳定的上行流是有限（proper）映射，wobbly 的则不是，但仍可以是支配的、重数非零的。[译者按（无外部依据）]}

我们从 Hitchin 截面上第 $k$ 个基本 Hecke 对应的推广概念开始。回忆\textbf{仿射 Grassmannian}
\[
\Gr:=\GL_n((z))/\GL_n[[z]],
\]
其中 $\GL_n((z)):=\GL_n(\C((z)))$，$\GL_n[[z]]:=\GL_n(\C[[z]])$。它是射影 ind 概形，特别地，其既约部分是射影簇的嵌套并。它参数化曲线一点处向量丛的高阶 Hecke 变换。

设
\[
\mu=(\mu_1\geq \mu_2\geq\dots \geq \mu_n)\in P^+\subset P\cong\Z^n
\]
是支配权，并设
\[
z^\mu:=\left(\begin{array}{cccc}z^{\mu_1} &0&\dots&0 \\ 0&z^{\mu_2}&\dots & 0\\ &&\vdots & \\ 0&0&\dots & z^{\mu_n}\end{array}\right)\in \GL_n(\C((z))).
\]
注意 $\GL_n[[z]]$ 从左边作用在 $\Gr$ 上，其轨道分解为
\[
\Gr=\coprod_{\mu\in P^+} \Gr^\mu=\coprod_{\mu\in P^+} \GL_n[[z]]\, [z^\mu],
\]
其中 $\Gr^\mu$ 以支配权 $\mu\in P^+$ 为标号，因为它们正是 $[z^\mu]$ 的轨道。有自然映射
\begin{equation}\label{fibr}
\Gr^\mu\to \GL_n [z^\mu] \cong \GL_n/P_\mu
\end{equation}
由令 $z=0$ 给出。注意 $\GL_n [z^\mu]\cong \GL_n/P_{\mu}$ 是部分旗簇，而映射 \eqref{fibr} 是 $\GL_n/P_{\mu}$ 上的一个有限秩向量丛。

用 $\oGr^\mu$ 记 $\Gr^\mu$ 在 $\Gr$ 中闭包的既约部分。于是有
\begin{equation}\label{partition}
\oGr^\mu=\coprod_{\mu \geq \lambda\in P^+}\Gr^\lambda,
\end{equation}
其中 $\mu \geq \lambda$ 是在 $P^+$ 上的\textbf{支配序}意义下，即 $\mu-\lambda$ 是若干正根之和（可带重数）。

一个重要例子是 $\mu=\omega_k=(\underbrace{1,\dots,1}_k,0,\dots,0)$ 为第 $k$ 个基本权的情形。此时 $\omega_k$ 是极小的（minuscule，即在支配序中极小），并且
\[
\oGr^{\omega_k}=\Gr^{\omega_k}\cong \Gr(k,n),
\]
即经典的 Grassmann 簇。

对支配权 $\mu$，现在可以在 Hitchin 截面上定义 $\mu$ 型的 Hecke 对应
\begin{equation}\label{general}
\calH^\mu:=\left\{ (\calE_a,V)\in W^+_0\times \Gr^\mu \;\Big|\;  \Phi_a|_{\Delta_c}(V)\subset V\right\}\subset W^+_0\times  \Gr^\mu,
\end{equation}
其中 $\Delta_c\to C$ 是 $c$ 在 $C$ 中的形式邻域。特别地 $\Delta_c\cong \Spec(\C[[z]])$。固定平凡化 $E_0|_{\Delta_c}\cong \calO^n_{\Delta_c}$ 与 $K|_{\Delta_c}\cong \calO_{\Delta_c}$，就可以把 $\gamma_a:=\Phi_a|_{\Delta_c}\in \gl_n[[z]]$ 看成矩阵。于是，若 $V=[g]\in \Gr^\mu$ 由 $g\in \GL_n((z))$ 代表，条件 $\Phi_a|_{\Delta_c}\in \gl_n[[z]]$ 就翻译为
\begin{equation}\label{affspring}
\gamma_a g^{-1}\in \gl_n[[z]].
\end{equation}
对固定的 $a\in \calA$，这定义了 $\gamma_a$ 的\textbf{仿射 Springer 纤维}。

在 $V\in \Gr^\mu$ 处对 $\calE_0$ 作 $\mu$ 型 Hecke 变换，由于不变性条件 \eqref{affspring}，将给出一个新的向量丛及其 Higgs 场。只要某些稳定性条件满足，新的 Higgs 丛就是稳定的，并且落在 $\calE_{\delta^\mu_c}$ 的上行流上；这里
\begin{equation}\label{deltamuc}
\delta^\mu_c=(\alpha_n c,\alpha_{n-1} c,\dots,\alpha_{1} c),
\end{equation}
而 $\mu=\sum_{i} \alpha_i \omega_i$ 是用基本权写出的，其中对 $1\leq i \leq n-1$ 有 $\alpha_i\in \N$，而 $\alpha_n\in \Z$。我们断言这样的 Hecke 变换诱导同构，于是得到下图：
\begin{equation}\label{diagram}
\begin{array}{ccc}
\calH^\mu&\stackrel{\cong}{\longrightarrow}& W^+_{\delta^\mu_c}\\
\pi_\mu{\Big\downarrow} && {\Big\downarrow} h_{{\delta^\mu_c}} \\
W^+_0 &\stackrel{h_0}{\longrightarrow}& \calA
\end{array}
\end{equation}
这样就有 $Q_{\pi_\mu}\cong Q_{h_{\delta^\mu_c}}$。于是，等变重数代数的计算被归结为 $\Gr^\mu$ 内部的计算，即研究描述 $\calH^\mu\subset W^+_0 \times \Gr^\mu$ 的方程。结果发现这归结为一个比较简单的线性代数计算。我们有下述结果与猜想。为了陈述它们，需要引入支配上行流及其重数的概念。

\begin{definition}
设 $\calE\in \M^{s\T}$。若 Hitchin 映射 $h_\calE:W^+_\calE\to \calA$ 是支配的，则称上行流 $W^+_\calE$ 是\textbf{支配的}。此时诱导的函数环映射 $h_\calE^*: \C[{\calA}]\to \C[{W^+_\calE}]$ 是单射，从而得到函数域扩张 $\C(\calA)\subset \C(W^+_\calE)$。定义该扩张的次数为 $W^+_\calE$ 的\textbf{重数}
\[
m_\calE:=[ \C(W^+_\calE):\C(\calA)].
\]
\end{definition}

注意 $m_\calE$ 也等于 $\C[\calA]$-模 $M_\calE$（由 $h_\calE^*$ 给出）的秩
\[
m_\calE=\dim_{\C(\calA)} (M_\calE\otimes_{\C[\calA]} \C(\calA)),
\]
也等于一般纤维的基数
\[
m_\calE=\#(h_\calE^{-1}(a)\cap W^+_\calE)
\]
（对一般的 $a\in\calA$）。注意后两种说法对一切上行流都有意义；而支配性等价于重数非零。

我们有下述结果：

\begin{theorem}[\cite{hausel-hitchin1}]\label{conjecturewobbly}
设 $\mu=(d+1)\omega_k\in P^+$，其中 $d\in \Z_{>0}$、$1\leq k\leq n-1$，并设 $c\in C$。当 $\calE_{\delta^\mu_c}\in \M^{s\T}$ 时，下述成立：
\begin{enumerate}
	\item $W^+_0$ 的 $\mu$ 型 Hecke 修改诱导同构 $\calH^\mu \to W^+_{\delta^\mu_c}$；
	\item $W^+_{\calE_{\delta^\mu_c}}$ 是支配的；
	\item $m_\calE=|W\cdot \mu|$，即 $\mu$ 的 Weyl 轨道的大小；
	\item 有
	\[
	Q_{\mu}\cong \C[J_d(Q_{\omega_k})]\cong \C\left[J_d(\Spec(H^*(\Gr(k,n),\C)))\right],
	\]
	其中对概形 $X$，我们用 $J_d(X)$ 记 $X$ 的 $d-1$ 阶 jet 概形。\footnote{译注：原文此处写作“the $d-1$th jet scheme of $\backslash$spec(R)”，其中 \texttt{\textbackslash spec} 在原文导言区并未定义（原文只有 \texttt{\textbackslash Spec}），直接在原 \texttt{.tex} 上编译会报 Undefined control sequence；由紧接的下一句 $J_d(X)(\C)=\Hom(\Spec(\C[z]/(z^d)),X)$ 判断，这里的意思是概形 $X$ 的 $d-1$ 阶 jet 概形，故译文按此表述。[来源：原文第 682 行及下一句的定义]}
	特别地，$J_d(X)(\C)=\Hom(\Spec(\C[z]/(z^d)),X)$ 是 $X$ 中 $d-1$ 阶 jet 的集合。
\end{enumerate}
\end{theorem}

\begin{remark}
我们猜想：(1)、(2)、(3) 对任意 $\mu \in P^+$ 都成立。
\end{remark}

\begin{remark}
作为 $n=2$ 情形的例子，我们给出 $d=1,2,3$ 时的重数代数。首先
\begin{equation}\label{d=1}
Q_{\delta_c^{2\omega_1}}\cong\C[a_0]/(a_0^2)\cong H^*(\P^1,\C).
\end{equation}
其次
\begin{equation}\label{d=2}
Q_{\delta_c^{2\omega_1}}\cong \C[a_0,a_1]/(a_0^2,a_0a_1).
\end{equation}
\footnote{译注：按定理 \ref{conjecturewobbly} 的标号，$\mu=(d+1)\omega_k$，故 $d=1,2,3$ 分别对应 $\mu=2\omega_1,3\omega_1,4\omega_1$。式 \eqref{d=1} 与 \eqref{d=2} 中的角标却都写作 $2\omega_1$，式 \eqref{d=3} 写作 $3\omega_1$；照上标号，\eqref{d=2} 与 \eqref{d=3} 处应分别作 $3\omega_1$ 与 $4\omega_1$。原文如此，疑为角标笔误，此处照录，读者按 $d=1,2,3$ 依次理解即可。[来源：原文式 (3.9)--(3.11) 与定理 3.12 的陈述]}
注意 $(a_0^2,a_0a_1)=(a_0,a_1)^2\cap (a_0)$，因此上行流 $W^+_{\delta_c^{2\omega_1}}$ 与幂零锥 $h^{-1}(0)$ 的概形论交 $\Spec(Q_{\delta_c^{2\omega_1}})$，是直线 $(a_0)$ 加上原点处一个二重嵌入点。注意这个上行流在 \cite[\S 8.2]{hausel-hitchin} 中研究过。

对 $d=3$，我们得到重数为 $2$ 的代数
\begin{equation}\label{d=3}
Q_{\delta_c^{3\omega_1}}\cong \C[a_0,a_1,a_2]/(a_0^2,a_0a_1,a_0a_2+a_1^2).
\end{equation}
\eqref{d=2} 与 \eqref{d=3} 都可由猜想 \ref{conjecturewobbly} 的第 (4) 条推出，二者也都可以按上面解释的办法在 $\Gr^\mu$ 中直接计算证明。
\end{remark}

\begin{remark}
$J_d(\Spec(H^*(\Gr(k,n),\C)))$ 竟然可以如此复杂，令人意外。特别地，在 $k=1$ 的情形（即射影空间上同调环的 jet 概形），关于其重数在 \cite[Conjecture III.21]{yuenthesis} 中只有一个猜想。
\end{remark}

\begin{remark}
最后我们指出：即使在 $(2)$ 型，也已经出现了新现象。如 \cite[\S 5.4]{hitchin4} 所讨论的，存在依赖于连续参数的重数代数；特别地，它们不可能同构于上同调环，因为上同调环是整的。
\end{remark}

\subsubsection{$W^+_\delta$ 的拉格朗日闭包}

\begin{definition}
设 $\calE\in \M^{s\T}$。$W^+_\calE$ 的\textbf{拉格朗日闭包} $\overline{\overline{W^+_\calE}}$ 是包含 $W^+_\calE$ 的最小闭的上行流之并。\footnote{译注：双重上划线是原文的记号，用来区别于通常的 Zariski 闭包；按原文的说明，它是按 Białynicki--Birula 分解取商之后所得的闭包，即“闭的上行流之并”里最小的那个。[来源：原文定义 3.17 同页的说明]}
换言之，拉格朗日闭包就是按 Białynicki--Birula 分解所取商空间中的闭包。
\end{definition}

利用 \eqref{diagram} 与 \eqref{partition}，可以推出下述结果。

\begin{theorem}[\cite{hausel-hitchin1}]
设 $\mu\in P^+$、$c\in C$，并由 \eqref{deltamuc} 回忆 $\delta^\mu_c$。假定 $\calE_{\delta^\mu_c}\in \M^{s\T}$。则
\[
\overline{\overline{W^+_\calE}}=\coprod_{\mu\geq \lambda\in P^+} W^+_{\delta^\lambda_c},
\]
即取遍支配序中不大于 $\mu$ 的支配权 $\lambda$ 所对应的上行流。
\end{theorem}

\subsection{走向几何 Satake 对应的经典极限}\label{geosat}


\origpage{2244}最后，我们来表述一些猜想，它们正是前面许多想法的原初动机。特别地，这些猜想暗示了 $\GL_n(\C)$ 的不可约表示的一种新构造，并更一般地暗示了任意复约化群 $\G$ 的不可约表示的新构造。

一般框架来自几何 Langlands 纲领的经典极限 \eqref{classlim}，其表述见 \cite{donagi-pantev}。这里我们用示意性的（并非完全严格的）方式，勾画这一经典极限的若干预期。它应当是某种凝聚层导出范畴之间的等价
\[
\mathscr{S}:D^b(\M_\Dol)\to D^b(\M_\Dol^\vee).
\]
这一等价的若干性质由 \cite{donagi-pantev} 提出，其中一部分已获证明。特别地，$\mathscr{S}$ 应当是沿 $\calA_\G\cong \calA_{\G^\vee}$ 中一般点的相对 Fourier--Mukai 变换。另一条关键性质 \cite{kapustin-witten}——即我们在上文 \S\ref{gls} 所称的增强镜像对称——是：$\mathscr{S}$ 应当把 $D^b(\M_\Dol)$ 上某些 Hecke 算子的作用与 $D^b(\M_\Dol)$ 上 Wilson 算子的作用相交换。\footnote{译注：此处原文两处都写作 $D^b(\M_\Dol)$，但按上文 \S\ref{gls} 中的定义，Wilson 算子 $\mathscr{W}^\mu$ 作用在镜像一侧 $D^b(\M^\vee_\Dol)$ 上（见下文 $\mathscr{W}^\mu$ 的定义）；式 \eqref{intertwine} 也表明了这一点。原文疑为笔误，此处照录，读者按 $\mathscr{S}$ 把 Hecke 算子送成 Wilson 算子来理解即可。[来源：原文 \S 3.4 首段与式 (3.12)]}

设 $\mu\in X_+(\G^\vee)=X^+(\G)$ 是 $\G^\vee$ 的支配特征。记
\[
\calH^\mu:=\{(E,\Phi)\in \M_\Dol, [g]\in \Gr^\mu \mid g^{-1}\Phi_c g \in \G[[z]]\}\subset \M_\Dol\times \Gr^\mu
\]
为某点 $c\in C$ 处的 Hecke 对应空间。它确实给出
\[
\begin{array}{ccc}
&\calH^\mu&\\
\pi_\mu\swarrow&&\searrow f^\mu \\
\M_\Dol&& \M_\Dol
\end{array}
\]
两个到 $\M_\Dol$ 的映射：第一个是到第一因子的投影 $\pi_\mu$，第二个是用相容的 Hecke 变换 $[g]\in \Gr^\mu$ 对 $(E,\Phi)$ 作 Hecke 变换\footnote{这里我们忽略稳定性方面的问题。}所得的 $f^\mu$。预期它诱导出
\[
\mathscr{H}^\mu :=f^\mu_* \circ \pi_\mu^*: D^b(\M_\Dol)\to  D^b(\M_\Dol),
\]
即 Hecke 算子，或称物理学家所说的 '\kern-0.1em t Hooft 算子。

另一方面，可以考虑所谓的 \textbf{Wilson 算子}
\[
\mathscr{W}^\mu : \begin{array}{ccc} D^b(\M^\vee_{Dol}) &\to& D^b(\M^\vee_\Dol) \\ \calF &\mapsto& \calF\otimes \rho_\mu(\bE|_{\M^\vee_{\Dol}\times \{c\}}) \end{array}
\]
它由张上万有 $\G^\vee$-丛 $\bE$ 在表示 $\rho_\mu: \G^\vee\to \GL(V_{\rho_\mu})$ 中所得的向量丛给出。

于是我们预期 \cite{donagi-pantev,kapustin-witten}
\begin{equation}\label{intertwine}
\mathscr{W}^\mu \circ \mathscr{S} = \mathscr{S} \circ \mathscr{H}^\mu.
\end{equation}

关于经典极限 $\mathscr{S}$ 还有两条预期，二者都受 Fourier--Mukai 变换（相应陈述在其中成立）的启发。第一，我们预期对任意 $\calF\in D^b(\M_\Dol)$ 有
\begin{equation}\label{Spush}
(h_\G)_*(\calF)\cong \mathscr{S}(\calF)|_{W^+_0}.
\end{equation}
第二，Hitchin 截面的结构层应当镜像于镜像 Higgs 模空间的结构层：
\begin{equation}\label{Sstructure}
\mathscr{S}(\calO_{W^+_0})\cong \calO_{\M^\vee_{\Dol}}.
\end{equation}

把 \eqref{intertwine} 与 \eqref{Sstructure} 结合起来，可以推出
\[
\mathscr{S}(\mathscr{H}^\mu(\calO_{W^+_0}))=\mathscr{W}^\mu(\mathscr{S}(\calO_{W^+_0}))=\mathscr{W}^\mu(\calO_{\M^\vee_{\Dol}}).
\]
一方面，我们应当有
\[
\mathrm{supp}(\mathscr{H}^\mu(\calO_{W^+_0}))=\overline{\overline{W^+_\mu}},
\]
其中 $W_\mu^+$ 是自某个 $\calE_\mu$ 出发的上行流，而 $\calE_\mu$ 是 $c\in C$ 处某 $\mu$ 型极大分裂 $\G$-Higgs 丛。另一方面，
\[
\mathscr{W}^\mu(\calO_{\M^\vee_{\Dol}})=\rho_\mu(\bE)_c=:\Lambda_\mu,
\]
即与主丛 $\bE_c$ 在表示 $\rho_\mu$ 中相伴的向量丛。

于是 Kapustin--Witten 的 \eqref{intertwine} 蕴含
\[
\mathscr{S}(\mathscr{H}^\mu(\calO_{W^+_0}))=\Lambda_\mu.
\]
用 \eqref{Spush} 可以检验这一点：
\[
\Lambda_\mu|_{ ^{L}W_0^+}= \mathscr{S}(\mathscr{H}^\mu(\calO_{W^+_0}))|_{ ^{L}W_0^+}=(h_\G)_*(\mathscr{H}^\mu(\calO_{W^+_0})).
\]

在 \cite{hausel-hitchin} 中我们论证了：非常稳定的 $(1,\dots,1)$ 型上行流 $W_{\delta}$ 的结构层，其镜像是
\[
\Lambda_\delta:=\bigotimes_{i=1}^{n-1} \bigotimes_{j=1}^{m_i} \Lambda^i\E_{c_{ij}},
\]
其中 $(\E,\bPhi)$ 是 $\M\times C$ 上的万有 Higgs 丛，而
\[
\delta_i=c_{i1}+c_{i2}+\dots \in C^{[m_i]}.
\]
特别地，镜像对称的一项预期是
\[
h_*(\calO_{W^+_\delta})\cong \Lambda_\delta|_{W^+_0}.
\]
这可由定理 \ref{mult111} 以及对 $\chi_\T (\E_c)$ 的直接计算得到。

在 \cite[\S 8.2]{hausel-hitchin} 中我们提出：当 $n=2$ 时，$Sym^2(\E_c)$ 的镜像应当是拉格朗日闭包 $\overline{\overline{W^+_{\delta_{c}^2}}}$ 的结构层，其中 $\delta_{c}^2=(0,2c)$。这可以推广如下。

\begin{conjecture}\label{geomsati}
设 $c\in C$，$\G$ 是约化群。则我们有下述猜想：
\begin{enumerate}
	\item 对任意 $\mu\in X_+(\G^\vee)$，$\rho_\mu(\bE_c)$ 的镜像的支集是 $\overline{\overline{W^+_{\delta_c^\mu}}}$；
	\item 设 $\mu\in X_+(\G^\vee)$ 使相应的不可约 $\G^\vee$-表示 $\rho_\mu$ 是重数自由的。则 $\rho_\mu(\bE_c)$ 的镜像是 $\calO_{\overline{\overline{W^+_{\delta_c^\mu}}}}$；
	\item 在后一情形，受限 Hitchin 映射 $h_\G:\overline{\overline{W^+_{\delta_c^\mu}}}\to \calA$ 的重数代数同构于 $\overline{\Gr}^\mu$ 的上同调环。
\end{enumerate}
\end{conjecture}

\begin{remark}
在 \cite{feigin-frenkel-rybnikov} 中，作者们研究几何 Langlands 纲领中的 opers，并对 $\G^\vee$ 的一切不可约表示 $\rho_\mu$ 的底向量空间 $V_\mu$ 构造了一个典范的 Poincaré 对偶环结构。当 $\rho_\mu$ 重数自由时，这个环同构于上同调环 $H^*(\overline{\Gr}^\mu)$。注意按 \cite[Theorem 1.5]{ginzburg}，这恰好就是
\[
H^*(\overline{\Gr}^\mu)\cong  I\!H^*(\overline{\Gr}^\mu),
\]
即上同调环满足 Poincaré 对偶的那些情形。此时这个环在 \cite{panyushev} 中得到了更细致的研究。
\end{remark}


\begin{center}
\rule{0.55\textwidth}{0.4pt}
\end{center}

\noindent\textbf{致谢。}\ 作者感谢 Nigel Hitchin 在 1995--1998 年间把他引向 Higgs 丛，在 1996 年提出 Langlands 对偶 Hitchin 系统的 SYZ 图景，以及近年来的合作 \cite{hausel-hitchin,hausel-hitchin1}。作者还感谢 David Ben-Zvi、Pierre-Henri Chaudouard、Pierre Deligne、Ron Donagi、Sergei Gukov、Jochen Heinloth、Joel Kamnitzer、G\'erard Laumon、Anton Mellit、David Nadler、Andy Neitzke、Ng\^{o} Bao Ch\^au、Michael Thaddeus、Tony Pantev、Du Pei、Rich\'ard Rim\'anyi、Leonid Rybnikov、Vivek Shende、Bal\'azs Szendr\H{o}i、Andr\'as Szenes、Fernando Rodriguez-Villegas、Edward Witten 与 Zhiwei Yun，多年来就本文所论主题所作的讨论。还要感谢 H\"ulya Arg\"uz、Jakub L\"owit、Bal\'azs Szendr\H{o}i 与 Nigel Hitchin 仔细阅读本文。
